\documentclass[../../main.tex]{subfiles}
\graphicspath{{../../image/}{../../figure/}}
\begin{document}

\section{紧性与相容性}\label{section:ch14-紧性与相容性}

本节将通过取 $\delta\downarrow0$ 极限来完成\cref{theorem:ch14-前沿追踪存在性}的证明。此处我们仅假设 $\{U_\delta\}$ 是\cref{definition:ch14-近似解}意义下任一 $\delta$-近似解族（$\delta$ 正且小），满足估计\eqref{align:ch14-delta-tv}与\eqref{align:ch14-delta-lip}。因而不需要前沿追踪算法所构造特定 $\delta$-近似解的特殊性质，例如激波以正确激波速度传播。

固定任一在 $(-\infty,\infty)\times[0,T)$ 中具有紧支集的检验函数 $\phi$。由 Green 定理，
\begin{align}
&\int_0^\infty\int_{-\infty}^\infty[\partial_t\phi U_\delta+\partial_x\phi F(U_\delta)]\,\mathrm{d}x\,\mathrm{d}t+\int_{-\infty}^\infty\phi(x,0)U_\delta(x,0)\,\mathrm{d}x\nonumber\\
&\qquad=-\int_0^\infty\sum\phi(y(t),t)\{F(U_\delta(y(t)+,t))-F(U_\delta(y(t)-,t))\nonumber\\
&\qquad\qquad\qquad\qquad\qquad-\dot y(t)[U_\delta(y(t)+,t)-U_\delta(y(t)-,t)]\}\,\mathrm{d}t,
\label{align:ch14-green}
\end{align}
其中对每个 $t$，求和遍历所有穿过 $t$-时间线的跳跃间断 $x=y(\cdot)$。

当跳跃间断 $x=y(\cdot)$ 是（近似）激波时，由\eqref{align:ch14-shock-speed}，
\begin{align}
&|F(U_\delta(y(t)+,t))-F(U_\delta(y(t)-,t))-\dot y(t)[U_\delta(y(t)+,t)-U_\delta(y(t)-,t)]|\nonumber\\
&\qquad\leqslant\delta|U_\delta(y(t)+,t)-U_\delta(y(t)-,t)|.
\label{align:ch14-residual-shock}
\end{align}
类似地，当 $x=y(\cdot)$ 是强度 $\leqslant\delta$ 的（近似）稀疏前沿时，由激波与稀疏波曲线间的接近性及\eqref{align:ch14-rarefaction-speed}，
\begin{align}
&|F(U_\delta(y(t)+,t))-F(U_\delta(y(t)-,t))-\dot y(t)[U_\delta(y(t)+,t)-U_\delta(y(t)-,t)]|\nonumber\\
&\qquad\leqslant c_5\delta|U_\delta(y(t)+,t)-U_\delta(y(t)-,t)|.
\label{align:ch14-residual-rarefaction}
\end{align}
最后，当 $x=y(\cdot)$ 是伪激波时，
\begin{align}
|F(U_\delta(y(t)+,t))-F(U_\delta(y(t)-,t))|\leqslant c_6|U_\delta(y(t)+,t)-U_\delta(y(t)-,t)|.
\label{align:ch14-residual-pseudo}
\end{align}

结合\eqref{align:ch14-residual-shock}、\eqref{align:ch14-residual-rarefaction}、\eqref{align:ch14-residual-pseudo}与\eqref{align:ch14-delta-tv}、\eqref{align:ch14-pseudoshock-strength}，我们推得对任一固定检验函数 $\phi$，\eqref{align:ch14-green}右端以 $C_\phi[TV_{(-\infty,\infty)}U_0(\cdot)+1]\delta$ 为界，因而当 $\delta\downarrow0$ 时趋于零。

由\eqref{align:ch14-delta-tv}、\eqref{align:ch14-delta-lip}与\cref{theorem:ch01-BV-三分}，任一收敛到零的 $\delta$ 序列都含子列 $\{\delta_k\}$，使 $\{U_{\delta_k}\}$ 在 $BV_{\mathrm{loc}}$ 中 a.e. 收敛到某 $U$。沿序列 $\{\delta_k\}$ 在\eqref{align:ch14-green}中取极限，并利用\eqref{align:ch14-initial-approx}，我们得到 $U$ 确实是\eqref{align:ch13-ivp}的弱解。



通过在\eqref{align:ch14-delta-tv}与\eqref{align:ch14-delta-lip}中取 $\delta\downarrow0$ 极限，可验证 $U$ 满足\eqref{align:ch13-step}与\eqref{align:ch13-strip}。此外，若 $\mathscr{C}$ 是任一关于所有 $U_\delta$ 非共振的 Lipschitz 曲线，则如\cref{section:ch14-总变差的界}所示，$U_\delta$ 在 $\mathscr{C}$ 上的迹有一致有界变差；于是取 $\delta\downarrow0$ 极限，我们推得 $U$ 在 $\mathscr{C}$ 上的迹有同样性质。



为完成证明，设 $(\eta,q)$ 是系统\eqref{align:ch14-system}的熵--熵流对，$\eta(U)$ 凸。令 $\phi$ 为 $(-\infty,\infty)\times[0,T)$ 上任一具有紧支集的非负检验函数。由 Green 定理，
\begin{align}
&\int_0^\infty\int_{-\infty}^\infty[\partial_t\phi\eta(U_\delta)+\partial_x\phi q(U_\delta)]\,\mathrm{d}x\,\mathrm{d}t+\int_{-\infty}^\infty\phi(x,0)\eta(U_\delta(x,0))\,\mathrm{d}x\nonumber\\
&\qquad=-\int_0^\infty\sum\phi(y(t),t)\{q(U_\delta(y(t)+,t))-q(U_\delta(y(t)-,t))\nonumber\\
&\qquad\qquad\qquad\qquad\qquad-\dot y(t)[\eta(U_\delta(y(t)+,t))-\eta(U_\delta(y(t)-,t))]\}\,\mathrm{d}t,
\label{align:ch14-entropy-green}
\end{align}
其中如\eqref{align:ch14-green}，对每个 $t$ 求和遍历所有穿过 $t$-时间线的跳跃间断 $x=y(\cdot)$。

当 $x=y(\cdot)$ 是（近似）激波时，熵不等式\eqref{align:ch08-熵条件}连同\eqref{align:ch14-shock-speed}蕴含
\begin{align}
&q(U_\delta(y(t)+,t))-q(U_\delta(y(t)-,t))-\dot y(t)[\eta(U_\delta(y(t)+,t))-\eta(U_\delta(y(t)-,t))]\nonumber\\
&\qquad\leqslant c_7\delta|U_\delta(y(t)+,t)-U_\delta(y(t)-,t)|.
\label{align:ch14-entropy-shock}
\end{align}
当 $x=y(\cdot)$ 是强度 $\leqslant\delta$ 的（近似）稀疏前沿时，\cref{theorem:ch08-弱激波熵产生}连同\eqref{align:ch14-rarefaction-speed}给出
\begin{align}
&|q(U_\delta(y(t)+,t))-q(U_\delta(y(t)-,t))-\dot y(t)[\eta(U_\delta(y(t)+,t))-\eta(U_\delta(y(t)-,t))]|\nonumber\\
&\qquad\leqslant c_8\delta|U_\delta(y(t)+,t)-U_\delta(y(t)-,t)|.
\label{align:ch14-entropy-rarefaction}
\end{align}
最后，当 $x=y(\cdot)$ 是伪激波时，
\begin{align}
&|q(U_\delta(y(t)+,t))-q(U_\delta(y(t)-,t))-\dot y(t)[\eta(U_\delta(y(t)+,t))-\eta(U_\delta(y(t)-,t))]|\nonumber\\
&\qquad\leqslant c_9|U_\delta(y(t)+,t)-U_\delta(y(t)-,t)|.
\label{align:ch14-entropy-pseudo}
\end{align}

结合\eqref{align:ch14-entropy-shock}、\eqref{align:ch14-entropy-rarefaction}、\eqref{align:ch14-entropy-pseudo}与\eqref{align:ch14-delta-tv}、\eqref{align:ch14-pseudoshock-strength}，我们推得对固定检验函数 $\phi$，\eqref{align:ch14-entropy-green}右端以 $-C_\phi[TV_{(-\infty,\infty)}U_0(\cdot)+1]\delta$ 为下界。因而沿序列 $\{\delta_k\}$ 取极限，我们得到解 $U$ 满足不等式\eqref{align:ch13-entropy-limit}，它表达熵可容性条件。\cref{theorem:ch14-前沿追踪存在性}的证明至此完成。



\end{document}
